\documentclass[aspectratio=169,t,xcolor=table]{beamer}
\usepackage{amsmath,amssymb}

% Tema Oficial do CEIA (Design System v3.0 - Research & Software Engineering Suite)
\usetheme{Ceia}

% Metadados Oficiais
\title{DESIGN SYSTEM CEIA v3.0}
\subtitle{SUÍTE VISUAL PARA PESQUISA CIENTÍFICA, ENGENHARIA DE SOFTWARE E IA}
\author{Departamento de Pesquisa \& Desenvolvimento CEIA}
\institute{Centro de Excelência em Inteligência Artificial -- UFG}
\date{Goiânia, 2026}

\begin{document}

% ==============================================================================
% 1. CAPA OFICIAL (Layout cover paramétrico TikZ)
% ==============================================================================
\setLayout{cover}
\begin{frame}[noframenumbering]
    \vspace{0.14\paperheight}
    \begin{columns}
        % Coluna Esquerda: Espaço ocupado pelo padrão oficial de gotas TikZ (Manual pág. 23)
        \column{0.44\textwidth}
        \phantom{Padrão Vetorial TikZ de Gotas}
        
        % Coluna Direita: Conteúdo e Logotipo Oficial
        \column{0.52\textwidth}
        {\fontsize{20}{24}\bfseries\color{white}\raggedright\inserttitle}\par\vspace{0.6em}
        {\fontsize{9.5}{13}\mdseries\color{CEIAGreen}\raggedright\insertsubtitle}\par\vspace{1.2em}
        
        \ceialogowhite[height=1.3cm]\par\vspace{1.2em}
        
        {\fontsize{8.5}{11}\mdseries\color{white!85}\insertauthor}\par\vspace{0.2em}
        {\fontsize{7.5}{9.5}\mdseries\color{white!65}\insertinstitute}
    \end{columns}
\end{frame}

% ==============================================================================
% 2. SUMÁRIO INSTITUCIONAL (\tableofcontents estilizado)
% ==============================================================================
\setLayout{content}
\begin{frame}{AGENDA DA APRESENTAÇÃO}{ESTRUTURAÇÃO DOS GRANDES CAPÍTULOS}
    \vspace{0.8em}
    \begin{center}
    \begin{minipage}{0.85\linewidth}
        \tableofcontents
    \end{minipage}
    \end{center}
\end{frame}

% ==============================================================================
% 3. ESTRUTURAÇÃO DE RACIOCÍNIO E LISTAS LIMPAS (Sem Caixas Fechadas)
% ==============================================================================
\setLayout{content}
\section{VISÃO GERAL \& DIRETRIZES}
\begin{frame}{ENGENHARIA DE SOFTWARE E IA}{ESTRUTURAÇÃO HIERÁRQUICA E DIRETRIZES TÉCNICAS}
    \vspace{0.1em}
    O \textbf{CEIA (Centro de Excelência em Inteligência Artificial)} atua na fronteira da pesquisa aplicada, desenvolvendo sistemas de IA com governança estrita e reprodutibilidade científica.
    \vspace{0.25em}

    \begin{columns}[T]
        \column{0.48\textwidth}
        {\small\bfseries\color{CEIADarkBlue} Pilares de Pesquisa \& Desenvolvimento}\par\vspace{0.15em}
        \begin{ceiaitemize}
            \item \textbf{Modelos Auditáveis}
            \begin{itemize}
                \item Soberania informacional e compliance
                \item Trilha de proveniência ponta a ponta
            \end{itemize}
            \item \textbf{Engenharia SDD}
            \begin{itemize}
                \item Contratos formais com Zod / Pydantic
            \end{itemize}
        \end{ceiaitemize}

        \column{0.48\textwidth}
        {\small\bfseries\color{CEIADarkBlue} Metodologia Operacional}\par\vspace{0.15em}
        \begin{ceiaitemize}
            \item \textbf{Supervisão Humana (HITL)}
            \begin{itemize}
                \item Validação de decisões críticas por experts
            \end{itemize}
            \item \textbf{Pipelines Determinísticos}
            \begin{itemize}
                \item Sandboxing estrito e retentativas bounded
                \item Tolerância a falhas e isolamento total
            \end{itemize}
        \end{ceiaitemize}
    \end{columns}
\end{frame}

% ==============================================================================
% 4. SEQUENCIAMENTO DE PROCESSOS (ceiastep)
% ==============================================================================
\begin{frame}{SEQUENCIAMENTO DE WORKFLOWS}{ETAPAS DO PIPELINE DE ENGENHARIA DE IA}
    \vspace{0.1em}
    \ceiastep{01}{Especificação Formal \& Contratos de Dados}{%
        Definição de interfaces inequívocas, schemas JSON e garantias criptográficas antes de qualquer escrita de código.
    }
    \ceiastep{02}{Execução por Agentes Especializados \& Linting}{%
        Delegação modular de tarefas de síntese para LLMs com validação sintática estrita via AST parser.
    }
    \ceiastep{03}{Gate Determinístico de Homologação (HITL)}{%
        Avaliação por gatekeeper humano assistido por testes automatizados de não-regressão e conformidade regulatória.
    }
\end{frame}

% ==============================================================================
% 5. COMPARAÇÃO E TRADE-OFFS (ceiavs & ceiaprogress)
% ==============================================================================
\begin{frame}{COMPARAÇÃO DE ABORDAGENS}{AVALIAÇÃO DE ROBUSTEZ: AD-HOC VS. ESPECIFICAÇÃO FORMAL}
    \ceiavs{Abordagem Legada (Ad-hoc)}{%
        \begin{itemize}
            \item Prompts soltos sem tipagem ou versionamento
            \item Extração frágil de dados baseada em regex
            \item Falhas silenciosas propagadas em runtime
        \end{itemize}
    }{Engenharia SDD (Padrão CEIA)}{%
        \begin{itemize}
            \item Contratos declarativos via Zod / Pydantic v2
            \item Structured Outputs com garantia de schema
            \item Trilha de auditoria criptográfica SHA-256
        \end{itemize}
    }
    \vspace{0.15em}
    \ceiaprogress{94}{Conformidade com Critérios de Aceite Automatizados}
    \ceiaprogress[CEIAGreen]{100}{Aderência Regulatória (LGPD \& EU AI Act Art. 12)}
\end{frame}

% ==============================================================================
% 6. RIGOR CIENTÍFICO E FORMULAÇÃO MATEMÁTICA (ceiatheorem & ceiaformula)
% ==============================================================================
\begin{frame}{RIGOR CIENTÍFICO E FORMULAÇÃO}{MODELAGEM TEÓRICA DE AMOSTRAGEM RESTRITA}
    \begin{ceiatheorem}[top=3.5pt, bottom=3.5pt]{Convergência Sob Restrições de Gramática Formal}
        Dada uma gramática livre de contexto $G = (V, \Sigma, R, S)$ e um modelo autorregressivo $P_\theta$, a amostragem mascarada restringe o suporte à linguagem $L(G)$ com probabilidade identicamente $1$, preservando a corretude sintática.
    \end{ceiatheorem}

    \begin{ceiaformula}[top=4pt, bottom=4pt]{Otimização por Entropia Cruzada com Penalização KL}
        \[
            \mathcal{L}_{\text{CEIA}}(\theta) = -\mathbb{E}_{(x, y) \sim \mathcal{D}} \left[ \sum_{t=1}^T \log P_\theta(y_t \mid y_{<t}, x) \right] + \lambda \, D_{\text{KL}}\left(P_\theta \,||\, P_{\text{ref}}\right)
        \]
        \vspace{1pt}
        {\scriptsize\color{CEIATextMuted}
        Onde $\theta$ representa os pesos treinados, $P_{\text{ref}}$ o modelo de referência e $\lambda$ a penalização KL.}
    \end{ceiaformula}
\end{frame}

% ==============================================================================
% 7. ARQUITETURA EM CAMADAS E LINHA DO TEMPO (ceiastack & ceiatimelinefour)
% ==============================================================================
\section{ARQUITETURA \& MULTI-AGENTE}
\ceiasectionframe{02}{ARQUITETURA \& MULTI-AGENTE}{Orquestração de Agentes, Contratos Formais e HITL Gates}

% ==============================================================================
% 7. ARQUITETURA EM CAMADAS E LINHA DO TEMPO (ceiastack & ceiatimelinefour)
% ==============================================================================
\begin{frame}{ARQUITETURA EM CAMADAS E CRONOGRAMA}{DECOMPOSIÇÃO DO SISTEMA E EVOLUÇÃO TEMPORAL}
    \vspace{0.1em}
    \ceialayer[CEIABrightBlue!8]{Camada 03: Orquestração e Agentes}{Workflows assíncronos via LangGraph com snapshotting e checkpointing}
    \ceialayer{Camada 02: Contratos e Governança}{JSON Schemas estritos, validadores Pydantic v2 e trilha auditável}
    \ceialayer{Camada 01: Execução e Isolamento}{Ambientes sandbox em contêineres e comunicação via gRPC seguro}
    
    \vspace{-0.1cm}
    \ceiatimelinefour{Q1 2026}{Contratos Zod}{Q2 2026}{AST Validator}{Q3 2026}{Multi-Agent}{Q4 2026}{Release v3.0}
\end{frame}

% ==============================================================================
% 8. CHECKLIST DE AUDITORIA E REQUISITOS (ceiachecklist)
% ==============================================================================
\begin{frame}{CHECKLIST DE AUDITORIA E REQUISITOS}{VERIFICAÇÃO DE PRONTIDÃO PARA PRODUÇÃO}
    \vspace{0.1em}
    \begin{ceiachecklist}
        \checkitem[ok]{Isolamento Criptográfico}{Hashes SHA-256 gerados para cada requisição arquivada.}
        \checkitem[ok]{Validação Sintática}{Schemas Zod com tipagem total dos payloads transitados.}
        \checkitem[warn]{Latência p99}{Fallback local ativado quando resposta externa $> 2.5\text{s}$.}
        \checkitem[ok]{Compilação Determinística}{Ambiente 100\% reprodutível via Docker e tectonic.}
    \end{ceiachecklist}

    \vspace{0.15em}
    \begin{ceiacard}[top=2pt, bottom=2pt]{Status de Homologação}
        {\footnotesize\color{CEIADarkBlue}\textbf{Parecer do Gatekeeper:} Homologado para produção com observabilidade contínua.}
    \end{ceiacard}
\end{frame}

% ==============================================================================
% 9. ESPECIFICAÇÕES FORMAIS SDD (Layout content)
% ==============================================================================
\begin{frame}{ENGENHARIA ORIENTADA A ESPECIFICAÇÕES}{GOVERNANÇA TÉCNICA E CONTRATOS DO SISTEMA}
    \vspace{0.2em}
    \ceiaspec{REQ-0101}{Isolamento Criptográfico e Trilha Auditável de Inferência}{CONCLUÍDO}
    \ceiaspec{REQ-0102}{Contratos de Dados com Schemas Zod e Saídas Estruturadas}{CONCLUÍDO}
    \ceiaspec{REQ-0103}{Fallback Determinístico com Retry Bounded para Chamadas de LLM}{CONCLUÍDO}
    \ceiaspec{REQ-0104}{Monitoramento de Drift e Conformidade LGPD \& EU AI Act}{EM EXECUÇÃO}
\end{frame}

% ==============================================================================
% 10. DIÁLOGO E PROMPTS DE AGENTES (Layout content)
% ==============================================================================
\begin{frame}{PROMPTS SISTÊMICOS E POLÍTICAS DE OPERAÇÃO}{ROLES FORMALIZADOS E DIRETRIZES DE SEGURANÇA}
    \vspace{0.2em}
    \begin{ceiaprompt}[ROLE: ORCHESTRATOR]{Diretrizes do Coordenador de Execução}
        Você é o orquestrador do pipeline de decisão autônoma. Decomponha o objetivo do usuário em etapas atômicas verificáveis. Nunca execute chamadas de ferramentas destrutivas sem aprovação explícita no gate humano G1.
    \end{ceiaprompt}
    \vspace{0.3em}
    \begin{ceiaprompt}[AGENT: SECURITY AUDITOR]{Verificação Contínua de Vulnerabilidades}
        Audite todos os payloads JSON antes do envio aos modelos locais. Se forem detectados padrões de injeção de prompt ou vazamento de credenciais, encerre a sessão com status BLOCKED e emita alerta de auditoria.
    \end{ceiaprompt}
\end{frame}

% ==============================================================================
% 11. JANELA DE CÓDIGO E DIFFS (Layout content com [fragile])
% ==============================================================================
\begin{frame}[fragile]{CONTRATOS DE DADOS E CODE DIFF}{REFACTORING E VALIDAÇÃO VIA ZOD}
\begin{ceiacode}[language=Java]{packages/core/src/agentContract.ts}
export const AgentVerdictSchema = z.object({
  specId: z.string(),
(*@\ceiaremoved{- approved: z.boolean(), // obsoleto: faltava trilha de auditoria}@*)
(*@\ceiaadded{+ status: z.enum(["APPROVED", "BLOCKED", "NEEDS\_REVISION"]),}@*)
(*@\ceiaadded{+ auditHash: z.string().min(32), // conformidade EU AI Act Art. 12}@*)
  issues: z.array(z.string()).default([]),
});
\end{ceiacode}
    \begin{ceiaalert}[top=1.5pt, bottom=1.5pt]{Structured Outputs com Verificação em Tempo de Execução}
        {\fontsize{7.2}{8.6}\selectfont Respostas de LLMs validadas contra schemas Zod. Saídas malformadas disparam retentativas determinísticas antes da propagação de estado.}
    \end{ceiaalert}
\end{frame}
\setLayout{content}

% ==============================================================================
% 12. PIPELINE MULTI-AGENTE EM TIKZ (Layout content)
% ==============================================================================
\begin{frame}{PIPELINE MULTI-AGENTE E GATES HITL}{GRAFO DE EXECUÇÃO DETERMINÍSTICO}
    \vspace{0.6em}
    \begin{center}
    \begin{tikzpicture}[node distance=0.5cm]
        \node[ceiaagent, minimum width=2.1cm] (disc) {Discovery\\Agent};
        \node[ceiagate, minimum width=1.6cm, right=of disc] (g1) {[G1]\\Gate};
        \node[ceiaagent, minimum width=2.1cm, right=of g1] (arch) {Software\\Architect};
        \node[ceiatool, minimum width=2.1cm, right=of arch] (tool) {Linter \&\\Zod Check};
        \node[ceiagate, minimum width=1.6cm, right=of tool] (g2) {[G2]\\Deploy};

        \draw[ceiaarr] (disc) -- (g1);
        \draw[ceiaarr] (g1) -- (arch);
        \draw[ceiaarr] (arch) -- (tool);
        \draw[ceiaarr] (tool) -- (g2);

        % Loop de feedback tracejado
        \draw[ceiafeedback] (tool.south) -- ++(0,-0.4) -| node[pos=0.25, below, font=\tiny\bfseries] {Rejeição / Refinamento} (arch.south);
    \end{tikzpicture}
    \end{center}
    \vspace{0.6em}
    \begin{ceiacard}{Semântica Visual do Grafo}
        \centering{\footnotesize \textbf{Azul Claro}: Agentes LLM \quad $\vert$ \quad \textbf{Verde Menta}: Ferramentas/Linters \quad $\vert$ \quad \textbf{Azul Escuro}: Gates de Decisão Humana}
    \end{ceiacard}
\end{frame}

% ==============================================================================
% 13. PIPELINE DECLARATIVO E ROADMAP (\ceiapipelinefour & \ceiaroadmapthree)
% ==============================================================================
\begin{frame}{PIPELINE DECLARATIVO E ROADMAP}{DEFINIÇÃO SEMÂNTICA DE FLUXO E CRONOGRAMA DE ENTREGAS}
    \vspace{0.4em}
    \ceiapipelinefour{Discovery\\Agent}{agent}{[G1] Decision\\Gate}{gate}{Software\\Architect}{agent}{Automated\\Deploy}{tool}
    \vspace{1.2em}
    \ceiaroadmapthree{FASE 1 (CONCLUÍDO)}{Elicitação \& Contratos}{FASE 2 (EM ANDAMENTO)}{Orquestração Local}{FASE 3 (PLANEJADO)}{Auditoria \& Escala}
\end{frame}

% ==============================================================================
% 14. TABELA DE BENCHMARK E EVAL DE MODELOS (Layout content)
% ==============================================================================
\section{BENCHMARK \& ESPECIFICAÇÕES}
\ceiasectionframe[bottom]{03}{BENCHMARKS \& ESPECIFICAÇÕES}{Avaliação Empírica de Modelos e Conformidade SDD}

\begin{frame}{AVALIAÇÃO DE MODELOS DE LINGUAGEM}{DESEMPENHO EM TAREFAS DE ENGENHARIA DE SOFTWARE}
    \vspace{0.1em}
    \setlength{\tabcolsep}{4.5pt}
    \begin{ceiatable}{l c c c r}
        \textbf{Modelo Avaliado} & \textbf{Pass@1 (\%)} & \textbf{Latência (s)} & \textbf{Custo / 1k Tok} & \textbf{Conformidade} \\
        \midrule
        Claude 3.5 Sonnet & \bestmetric{92.4\%} & 1.42s & \$0.003 & 100\% \\
        GPT-4o & 90.2\% & \bestmetric{1.10s} & \$0.005 & 98\% \\
        Llama-3-70B (Local) & 82.6\% & 2.15s & \bestmetric{\$0.000} & 100\% (LGPD) \\
        DeepSeek-Coder-V2 & 88.7\% & 1.65s & \$0.001 & 95\% \\
    \end{ceiatable}
    \vspace{0.15em}
    \begin{ceiablock}[top=3.5pt, bottom=3.5pt]{Critério de Seleção para Ambientes de Produção}
        {\scriptsize O modelo Sonnet alcança o maior índice de aprovação sintática sem intervenção, enquanto modelos locais garantem isolamento total de dados sensíveis.}
    \end{ceiablock}
\end{frame}

% ==============================================================================
% 15. MATRIZ EXECUTIVA E ÍCONES VETORIAIS (ceiamatrixtable & \ceiaicon)
% ==============================================================================
\begin{frame}{MATRIZ COMPARATIVA DE CAPACIDADES}{ANÁLISE DE CONFORMIDADE E ÍCONES VETORIAIS TIKZ}
    \vspace{0.1em}
    \begin{ceiamatrixtable}{l c c c}
        \textbf{Capacidade Avaliada} & \textbf{Abordagem Legada} & \textbf{Cloud Proprietária} & \textbf{Padrão CEIA} \\
        \midrule
        \ceiaicon{shield}\hspace{4pt} Governança \& LGPD & \matrixfail & \matrixwarn & \matrixok (Local) \\
        \ceiaicon{code}\hspace{4pt} Contratos Zod Tipados & \matrixfail & \matrixfail & \matrixok (Formal) \\
        \ceiaicon{server}\hspace{4pt} Soberania de Dados & \matrixwarn & \matrixfail & \matrixok (100\%) \\
        \ceiaicon{clock}\hspace{4pt} Latência de Inferência & 3.2s & 1.8s & \bestmetric{0.45s} \\
    \end{ceiamatrixtable}
    \vspace{0.1em}
    \ceiabanner[ok]{Conformidade Legal Assegurada}{Todas as métricas atendem integralmente ao Art. 12 do EU AI Act.}
\end{frame}

% ==============================================================================
% 16. DIRETRIZES DE ENGENHARIA E BANNERS (ceiadatatable & \ceiabanner)
% ==============================================================================
\begin{frame}{ESPECIFICAÇÕES TÉCNICAS E DIRETRIZES}{TABELA COM AUTO-WRAP E BANNERS DE SÍNTESE}
    \begin{ceiadatatable}{l c p{8.2cm}}
        \textbf{Módulo} & \textbf{Status} & \textbf{Diretriz Operacional do Sistema} \\
        \midrule
        \ceiaicon{brain}\hspace{3pt} Runtime LLM & \bestmetric{ATIVO} & Execução local com modelos quantizados em 4-bit para garantia de isolamento. \\
        \ceiaicon{database}\hspace{3pt} Base Vetorial & \bestmetric{ATIVO} & Indexação HNSW de embeddings criptografados com rotação semanal de chaves. \\
    \end{ceiadatatable}
    \vspace{0.3em}
    \ceiabanner[warn]{Atenção}{Modelos proprietários externos não devem receber dados protegidos pela LGPD.}
    \ceiabanner[info]{Tese Central}{Arquiteturas baseadas em contratos reduzem o tempo de depuração em até 70\%.}
\end{frame}

% ==============================================================================
% 17. GRÁFICOS MOLDÁVEIS: DISTRIBUIÇÃO E COMPARAÇÃO (ceiadonut & ceiavbars)
% ==============================================================================
\begin{frame}{GRÁFICOS VETORIAIS MOLDÁVEIS}{DISTRIBUIÇÃO DE RECURSOS E HISTOGRAMA COMPARATIVO}
    \begin{columns}[c]
        \column{0.48\textwidth}
        \begin{ceiacard}[top=2pt, bottom=2pt]{Distribuição de Inferência}
            \centering\vspace{0.1em}
            \ceiadonutthree{55}{CEIADarkBlue}{30}{CEIABrightBlue}{15}{CEIAGreen}{3.2M\\Reqs}\par\vspace{0.2em}
            {\tiny \textbf{Azul Escuro}: Local (55\%) \quad \textbf{Azul}: Fine-tuned (30\%) \quad \textbf{Verde}: Cloud (15\%)}
        \end{ceiacard}

        \column{0.48\textwidth}
        \begin{ceiacard}[top=2pt, bottom=2pt]{Acurácia por Benchmark (\%)}
            \centering\vspace{0.1em}
            \ceiavbarsthree{88}{HumanEval}{94}{MBPP}{91}{GSM8K}\par\vspace{0.2em}
            {\tiny Taxa de acerto superior a 85\% em todos os domínios avaliados.}
        \end{ceiacard}
    \end{columns}
\end{frame}

% ==============================================================================
% 18. COMPOSIÇÃO DE FIGURA LATERAL (ceiasidefigure)
% ==============================================================================
\section{IDENTIDADE \& IMPACTO}
\begin{frame}{COMPOSIÇÃO DE FIGURA LATERAL}{LAYOUT CANÔNICO: DIAGRAMA À ESQUERDA E EXPLICAÇÃO À DIREITA}
    \begin{ceiasidefigure}{lib/logos/ceia_simbolo.pdf}{Símbolo Oficial CEIA}
        {\small\bfseries\color{CEIADarkBlue} Princípios de Identidade Visual}\par\vspace{0.3em}
        \begin{itemize}
            \setlength{\itemsep}{0.3em}
            \item \textbf{Cérebro Conectado}: Sinergia entre pesquisa e ecossistema produtivo.
            \item \textbf{Geometria Oficial}: Arcos com raio $R = W/2$ e ângulo reto de $90^\circ$.
            \item \textbf{Área de Proteção}: Distância mínima de $2x$ mantida em todas as aplicações.
        \end{itemize}
    \end{ceiasidefigure}
\end{frame}

% ==============================================================================
% 19. INDICADORES DE IMPACTO DO SISTEMA (Layout content)
% ==============================================================================
\begin{frame}{INDICADORES DE IMPACTO DO SISTEMA}{GANHOS COMPROVADOS EM AMBIENTE DE PRODUÇÃO}
    \vspace{1.2em}
    \ceiastatfour{99.1\%}{F1-Score Geral}{4.2x}{Throughput}{18ms}{Inferência}{0}{Falhas Críticas}
\end{frame}

% ==============================================================================
% 20. SUBMARCAS INSTITUCIONAIS (Layout content)
% ==============================================================================
\setLayout{soft}
\begin{frame}{ECOSSISTEMA E EXTENSÕES DE MARCA}{SUBMARCAS OFICIAIS CONFORME O MANUAL DE IDENTIDADE}
    \begin{columns}[c]
        \column{0.48\textwidth}
        \begin{ceiacard}[top=3pt, bottom=3pt]{Unidade Educacional}
            \centering\vspace{0.3em}
            \ceiasubmarca{Academy}\par\vspace{0.4em}
            {\scriptsize Formação técnica avançada em IA e engenharia de software.}
        \end{ceiacard}

        \column{0.48\textwidth}
        \begin{ceiacard}[top=3pt, bottom=3pt]{Internacionalização}
            \centering\vspace{0.3em}
            \ceiasubmarca{Global}\par\vspace{0.4em}
            {\scriptsize Cooperação científica global e atração de talentos de ponta.}
        \end{ceiacard}
    \end{columns}
    \par\vspace{0.2em}
    \begin{ceiacard}[top=2pt, bottom=2pt]{Régua Institucional de Parceiros}
        \centering\vspace{0.1em}
        \ceiaregua[width=0.75\textwidth]\par
    \end{ceiacard}
\end{frame}
\setLayout{content}

% ==============================================================================
% 21. CITAÇÃO ACADÊMICA E PILARES EXECUTIVOS (ceiaquote & ceiapillars)
% ==============================================================================
\begin{frame}{FUNDAMENTAÇÃO E PILARES ESTRATÉGICOS}{DIRETRIZES CONSOLIDADAS DO CENTRO DE EXCELÊNCIA}
    \begin{ceiaquote}[Vaswani et al., NeurIPS 2017]{Attention Is All You Need}
        The dominant sequence transduction models are based on complex recurrent or convolutional neural networks that include an encoder and a decoder.
    \end{ceiaquote}
    \vspace{0.2em}
    \ceiapillars{1}{Pesquisa Aplicada}{Desenvolvimento de algoritmos estado-da-arte com publicação de impacto.}{2}{Governança \& Ética}{Trilhas de auditoria criptográfica e total conformidade com a LGPD.}{3}{Engenharia de IA}{Sistemas escaláveis, contratos formais e observabilidade ponta a ponta.}
\end{frame}

% ==============================================================================
% 22. SUMARIZAÇÃO HERO / PUNCHY (Pouco Texto: 1 Grande Conclusão de Alto Impacto)
% ==============================================================================
\section{SÍNTESE \& SUMARIZAÇÃO}
\begin{frame}{SUMARIZAÇÃO HERO: CONCLUSÃO DE ALTO IMPACTO}{MODELO PARA POUCO TEXTO COM MÁXIMA ENERGIA VISUAL}
    \vspace{0.3em}
    \begin{ceiasummarypunch}{Veredito de Engenharia \& Pesquisa}{Arquitetura Homologada para Operação Crítica em Produção}
        A migração para orquestração declarativa e contratos formais de tipos garantiu estabilidade determinística, eliminando 100\% dos incidentes de parsing silencioso.
        \par\vspace{1.2em}
        \noindent
        \ceiachip[ok]{shield}{Conformidade LGPD \& EU AI Act}\hfill
        \ceiachip[info]{server}{Soberania On-Premise}\hfill
        \ceiachip[warn]{clock}{Latência p95 < 120ms}
    \end{ceiasummarypunch}
\end{frame}

% ==============================================================================
% 23. SUMARIZAÇÃO SPLIT (Texto Médio: Veredito Executivo + Evidências)
% ==============================================================================
\begin{frame}{SUMARIZAÇÃO SPLIT: VEREDITO E EVIDÊNCIAS}{MODELO ASSIMÉTRICO: CARD DE IMPACTO À ESQUERDA E DESDOBRAMENTOS À DIREITA}
    \vspace{0.1em}
    \begin{ceiasummarysplit}[0.38]{Veredito Final}{99.4\%}{Sistema aprovado em todos os testes de estresse e homologação técnica.}{shield}
        \ceiasplititem[CEIABrightBlue]{Orquestração Determinística}{Grafo multi-agente executa com fallback automático para modelos locais.}
        \ceiasplititem[CEIAGreen]{Eficiência de Tokens \& Custo}{Redução comprovada de 68\% no consumo de tokens via caching semântico.}
        \ceiasplititem[CEIADarkBlue]{Observabilidade Contínua}{Trilhas OpenTelemetry com rastreamento fim a fim de cada chamada LLM.}
    \end{ceiasummarysplit}
\end{frame}

% ==============================================================================
\begin{frame}{SUMARIZAÇÃO COMPACTA: SÍNTESE MULTI-PONTO}{MODELO BALANCEADO PARA GRANDE VOLUME DE TEXTO SEM ESTOURAR O SLIDE}
    \begin{columns}[T]
        \begin{column}{0.485\textwidth}
            \ceiacompactitem{\matrixok}{Contratos Tipados}{Interfaces Zod validadas estritamente antes de qualquer inferência.}
            \ceiacompactitem{\matrixok}{HITL Gates Operacionais}{Aprovação humana mandatória antes de deploys em produção.}
            \ceiacompactitem{\matrixok}{Cobertura de Testes 94\%}{Garantia de regressão formal sobre pipelines multi-agente.}
        \end{column}
        \begin{column}{0.485\textwidth}
            \ceiacompactitem{\matrixok}{Latência Sob SLA}{Tempo de inferência p95 mantido abaixo de 180ms em regime de carga.}
            \ceiacompactitem{\matrixwarn}{Cache Semântico}{Em homologação final com ganho adicional projetado de 20\% em tokens.}
            \ceiacompactitem{\matrixok}{Soberania On-Premise}{Zero tráfego de dados confidenciais para clouds externas.}
        \end{column}
    \end{columns}
    \vspace{0.1em}
    \ceiabanner[info]{Conclusão da Auditoria}{Todos os requisitos obrigatórios foram atendidos integralmente pelo projeto.}
\end{frame}

% ==============================================================================
% 16. TESE INSTITUCIONAL COM ACTAY WIDE (Layout statement)
% ==============================================================================
\setLayout{statement}
\begin{frame}
    \vspace{0.12\paperheight}
    \begin{center}
        \ceiaweareceia[height=2.2cm]\par\vspace{1.2em}
        {\fontsize{15}{19}\mdseries\color{CEIABlue}O FUTURO DA TECNOLOGIA JÁ TEM ENDEREÇO.}\par\vspace{1.5em}
        {\small\color{CEIATextMuted}CENTRO DE EXCELÊNCIA EM INTELIGÊNCIA ARTIFICIAL -- UFG}
    \end{center}
\end{frame}

% ==============================================================================
% 17. ENCERRAMENTO INSTITUCIONAL (Layout closing oficial)
% ==============================================================================
\setLayout{closing}
\begin{frame}[noframenumbering]
    \begin{tikzpicture}[remember picture, overlay]
        % Centro absoluto: OBRIGADO :)
        \node[anchor=center, align=center] at ($(current page.center)+(0, 0.02\paperheight)$) {
            {\fontsize{44}{50}\selectfont\bfseries\color{white} OBRIGADO :)}
        };

        % Rodapé centralizado e perfeitamente equilibrado (Manual pág. 33)
        \node[anchor=south, align=center] at ($(current page.south)+(0, 0.75)$) {
            {\fontsize{11}{14}\selectfont\color{CEIAGreen}\bfseries ceia.ufg.br \quad\textbar\quad contato@ceia.ufg.br}\\[6pt]
            {\fontsize{8}{10}\selectfont\color{white!70} CENTRO DE EXCELÊNCIA EM INTELIGÊNCIA ARTIFICIAL --- UFG}
        };
    \end{tikzpicture}
\end{frame}

\end{document}
